\documentclass[12pt]{article}

\usepackage[utf8]{inputenc}
\usepackage[T1]{fontenc}
\usepackage{lmodern}
\usepackage[spanish]{babel}
\usepackage{booktabs}
\usepackage{array}
\usepackage{tabularx}
\usepackage{longtable}
\usepackage{enumitem}
\usepackage[hidelinks]{hyperref}
\usepackage{xurl}

\sloppy
\setlength{\parindent}{0pt}

\begin{document}

\begin{center}
  {\LARGE \textbf{Trabajo Práctico Integrador}}\\[0.5em]
  {Analítica de Datos, Universidad de San Andrés}
\end{center}

\section{Generalidades}

El trabajo es grupal e integra los contenidos de la materia: desde EDA y preparación de datos hasta entrenamiento y evaluación de modelos de Machine Learning supervisado.

\begin{itemize}
    \item Grupos de \textbf{3 o 4 personas}.
    \item Un mismo dataset no puede ser elegido por más de dos grupos.
    \item Inscripción del grupo y del dataset elegido en la planilla compartida.
    \item \textbf{Fecha límite de inscripción: jueves 21/05 (clase 9).}
\end{itemize}

El trabajo se entrega y se defiende en \textbf{dos tandas}. Después de la primera defensa los grupos reciben una devolución que deben incorporar antes de la segunda entrega.

\section{Datasets}

Pueden elegir uno de la lista o proponer uno propio (sujeto a aprobación antes de la fecha límite). Para cada uno se indica el tema y el tipo de problema sugerido.

Datasets pesados (Yellow Cab, HVFHS, Yelp completo) requieren acotar el alcance y justificarlo. Datasets sensibles (policiales, judiciales, salud) requieren una sección de sesgos y consideraciones éticas.

\subsection{AirBnB Buenos Aires}
\begin{itemize}
    \item \textbf{Tema:} alojamiento.
    \item \textbf{Tipo sugerido:} regresión / clasificación.
    \item \textbf{Link:} \url{https://insideairbnb.com/buenos-aires/}
\end{itemize}

\subsection{Encuesta Mundial de Salud Escolar (EMSE 2018)}
\begin{itemize}
    \item \textbf{Tema:} educación / salud.
    \item \textbf{Tipo sugerido:} clasificación.
    \item \textbf{Link:} \url{http://datos.salud.gob.ar/dataset/base-de-datos-de-la-3-encuesta-mundial-de-salud-escolar-emse-con-resultados-nacionales-argentina/archivo/9888b372-8094-45e2-8b62-40630838c5dc}
\end{itemize}

\subsection{Estadísticas sobre Ejecución de la Pena (SNEEP)}
\begin{itemize}
    \item \textbf{Tema:} judicial.
    \item \textbf{Tipo sugerido:} clasificación. Elegir un año.
    \item \textbf{Link:} \url{https://datos.jus.gob.ar/dataset/sneep/archivo/5fd7ce53-c741-4837-9850-d2879fec8a6b}
\end{itemize}

\subsection{Recorridos en Ecobicis}
\begin{itemize}
    \item \textbf{Tema:} tránsito.
    \item \textbf{Tipo sugerido:} regresión / clasificación. Elegir un año.
    \item \textbf{Link:} \url{https://data.buenosaires.gob.ar/dataset/bicicletas-publicas}
\end{itemize}

\subsection{Crímenes reportados en Chicago}
\begin{itemize}
    \item \textbf{Tema:} policial.
    \item \textbf{Tipo sugerido:} clasificación. Elegir un año.
    \item \textbf{Link:} \url{https://data.cityofchicago.org/Public-Safety/Crimes-2001-to-Present/ijzp-q8t2/related_content}
\end{itemize}

\subsection{Denuncias a la policía de NY}
\begin{itemize}
    \item \textbf{Tema:} policial.
    \item \textbf{Tipo sugerido:} clasificación. Datos 2024--2025.
    \item \textbf{Link:} \url{https://data.cityofnewyork.us/Public-Safety/NYPD-Complaint-Data-Current-Year-To-Date-/5uac-w243/about_data}
\end{itemize}

\subsection{Yellow Cab NYC}
\begin{itemize}
    \item \textbf{Tema:} tránsito.
    \item \textbf{Tipo sugerido:} regresión. Acotar mes/año por el tamaño.
    \item \textbf{Link:} \url{https://www1.nyc.gov/site/tlc/about/tlc-trip-record-data.page}
\end{itemize}

\subsection{Eventos de violencia organizada (UCDP GED)}
\begin{itemize}
    \item \textbf{Tema:} conflictos bélicos.
    \item \textbf{Tipo sugerido:} clasificación.
    \item \textbf{Link:} \url{https://ucdp.uu.se/downloads/index.html#ged_global}
\end{itemize}

\subsection{TMDB Movies Dataset 2024}
\begin{itemize}
    \item \textbf{Tema:} ocio.
    \item \textbf{Tipo sugerido:} regresión.
    \item \textbf{Link:} \url{https://www.kaggle.com/datasets/asaniczka/tmdb-movies-dataset-2023-930k-movies}
\end{itemize}

\subsection{Pacientes con diabetes en 130 hospitales (US)}
\begin{itemize}
    \item \textbf{Tema:} salud.
    \item \textbf{Tipo sugerido:} clasificación (readmisión).
    \item \textbf{Link:} \url{https://archive.ics.uci.edu/dataset/296/diabetes+130-us+hospitals+for+years+1999-2008}
\end{itemize}

\subsection{Precios Claros --- Base SEPA}
\begin{itemize}
    \item \textbf{Tema:} comercio.
    \item \textbf{Tipo sugerido:} regresión / clasificación. Elegir una cadena.
    \item \textbf{Link:} \url{https://datos.gob.ar/dataset/produccion-precios-claros---base-sepa}
\end{itemize}

\subsection{High-Volume For-Hire Services (HVFHS)}
\begin{itemize}
    \item \textbf{Tema:} tránsito.
    \item \textbf{Tipo sugerido:} regresión / clasificación. Elegir un año.
    \item \textbf{Link:} \url{https://www1.nyc.gov/site/tlc/about/tlc-trip-record-data.page}
\end{itemize}

\subsection{Reviews de Yelp}
\begin{itemize}
    \item \textbf{Tema:} comercio.
    \item \textbf{Tipo sugerido:} clasificación.
    \item \textbf{Link:} \url{https://business.yelp.com/data/resources/open-dataset/}
\end{itemize}

\subsection{Dataset propio}
\begin{itemize}
    \item \textbf{Tema:} a definir por el grupo.
    \item \textbf{Tipo sugerido:} a justificar.
    \item \textbf{Link:} sujeto a aprobación previa.
\end{itemize}

\section{Consignas}

\subsection{Entrega 1 --- EDA y preparación de datos}

\subsubsection{EDA}

\begin{itemize}
    \item Cargar el dataset y describir origen, contexto y forma de recolección.
    \item Número de observaciones, número de variables, tipos de datos.
    \item Distribuciones marginales y relaciones bivariadas.
    \item Errores, outliers y valores faltantes. Clasificar los faltantes (MCAR, MAR, MNAR).
\end{itemize}

\subsubsection{Visualización}

\begin{itemize}
    \item Elegir el gráfico adecuado según el tipo de variable y el análisis.
    \item Justificar la elección de cada gráfico.
    \item Interpretar lo que se ve.
\end{itemize}

\subsubsection{Definición del problema de ML supervisado}

\begin{itemize}
    \item Plantear un problema de clasificación o regresión.
    \item Definir la variable target.
    \item Definir las métricas de evaluación y justificar su elección.
\end{itemize}

\subsubsection{Preprocesamiento}

\begin{itemize}
    \item Limpieza general del dataset.
    \item Split (train/test o train/validation/test). Justificar la estrategia.
    \item Tratamiento de valores faltantes.
    \item Detección y manejo de outliers.
    \item Escalado/normalización cuando corresponda.
\end{itemize}

El preprocesamiento que use estadísticas del dataset debe ajustarse solo con train.

\subsubsection{Feature engineering}

\begin{itemize}
    \item Creación de nuevos features justificada.
    \item Codificación de variables categóricas.
    \item Discretización cuando corresponda.
    \item En clasificación: análisis de balance de clases y, si corresponde, técnica de balanceo justificada.
\end{itemize}

\subsubsection{Reducción de dimensionalidad}

\begin{itemize}
    \item Selección de features con al menos un método (filtros, correlación, importancia desde un modelo base, etc.).
    \item PCA: implementación, varianza explicada y discusión.
\end{itemize}

\subsection{Entrega 2 --- Modelado y conclusiones}

\subsubsection{Modelado}

\begin{itemize}
    \item Un baseline simple (regresión lineal/logística, árbol o dummy).
    \item Al menos dos modelos más vistos en clase (KNN, Random Forest, Boosting, Naive Bayes, etc.).
    \item Validación cruzada sobre train.
    \item Tuning de hiperparámetros del mejor modelo.
    \item Tabla comparativa de resultados.
    \item Performance final sobre test, reportada una sola vez.
\end{itemize}

\subsubsection{Interpretación y conclusiones}

\begin{itemize}
    \item Interpretabilidad: feature importance, coeficientes o herramientas similares (SHAP, permutation importance).
    \item Conclusiones sobre el uso del modelo y sus limitaciones.
    \item Sesgos del dataset y consideraciones éticas (obligatorio en datasets sensibles).
    \item Incorporación de la devolución recibida en la Entrega 1.
\end{itemize}

\section{Entregas}

Cada entrega incluye:

\begin{enumerate}
    \item \textbf{Repositorio de GitHub} actualizado con el código de la etapa correspondiente. \texttt{README.md} y \texttt{requirements.txt} incluidos. No subir el dataset crudo si es pesado.
    \item \textbf{Informe parcial}: notebook ejecutado de punta a punta, con celdas de markdown explicando cada paso.
    \item \textbf{Presentación}: hasta 6 diapositivas + carátula por entrega, con link al repositorio.
\end{enumerate}

\begin{center}
\begin{tabular}{l l l}
\toprule
\textbf{Entrega} & \textbf{Contenido} & \textbf{Fecha límite} \\
\midrule
Entrega 1 & EDA y preparación de datos & martes 09/06 \\
Entrega 2 & Modelado y conclusiones & martes 23/06 \\
\bottomrule
\end{tabular}
\end{center}

\section{Defensas orales}

Cada entrega se defiende oralmente en clase. La devolución de la primera defensa debe incorporarse en la segunda entrega.

\begin{center}
\begin{tabular}{l l l}
\toprule
\textbf{Defensa} & \textbf{Clase} & \textbf{Fecha} \\
\midrule
Defensa 1 --- EDA y preparación & Clase 13 & 18/06 \\
Defensa 2 --- Modelado y conclusiones & Clase 14 & 25/06 \\
\bottomrule
\end{tabular}
\end{center}

\begin{itemize}
    \item 15 minutos por grupo: 10' de exposición + 5' de preguntas.
    \item Todos los integrantes exponen una parte aproximadamente igual.
    \item Cualquier integrante puede recibir preguntas sobre cualquier parte del trabajo.
\end{itemize}

\section{Evaluación}

La nota final se compone de las dos entregas, con peso aproximadamente igual.

\begin{center}
\begin{tabular}{l c c}
\toprule
\textbf{Dimensión} & \textbf{Peso} & \textbf{Entrega} \\
\midrule
Entendimiento del dataset y del problema & 5\% & 1 \\
EDA y visualizaciones & 8\% & 1 \\
Preprocesamiento y feature engineering & 8\% & 1 \\
Reducción de dimensionalidad & 5\% & 1 \\
Modelado: elección, validación, comparación, tuning & 10\% & 2 \\
Interpretación y conclusiones & 5\% & 2 \\
Incorporación de la devolución de la Entrega 1 & 3\% & 2 \\
Repositorio y reproducibilidad & 6\% & 1 y 2 \\
Presentación y defensa oral & 50\% & 1 y 2 \\
\bottomrule
\end{tabular}
\end{center}

No se aceptan trabajos sin componente de modelado.

\section{Uso de inteligencia artificial}

Se incentiva el uso de herramientas de inteligencia artificial (ChatGPT, Claude, Copilot, etc.) durante todo el trabajo práctico: para explorar el dataset, generar código, debuggear, redactar el informe o lo que necesiten. La única condición es que entiendan perfecto qué carajo están haciendo. Si en la defensa oral no saben contestar algo porque lo hizo la IA y no ustedes, se penalizará fuerte.

\end{document}
